\chapter{系统设计}

\section{总体架构}

系统总体流程如图像采集、目标检测、状态机、控制器和舵机驱动五个部分组成。摄像头采集到的画面首先进入视觉检测模块，输出目标是否存在、目标中心、候选框和面积等信息；状态机根据连续检测结果判断当前处于等待、跟踪、短时丢失或搜索状态；控制模块根据目标中心偏差计算下一帧云台角度；硬件驱动模块通过 PCA9685 输出 PWM 信号控制两个舵机。

\begin{figure}[htbp]
  \centering
  \fbox{\parbox{0.86\linewidth}{
  \centering
  摄像头采集 $\rightarrow$ HSV 目标检测 $\rightarrow$ 状态机判断 $\rightarrow$ PID/限位控制 $\rightarrow$ PCA9685 $\rightarrow$ 双舵机云台
  }}
  \caption{系统总体闭环流程}
  \label{fig:system-flow}
\end{figure}

\section{硬件组成}

硬件部分包括 OrangePi 5 Pro、USB 摄像头、PCA9685 16 路 PWM 舵机驱动板、两个舵机和两自由度云台结构。当前已验证的舵机通道配置为：PCA9685 的 channel 0 对应俯仰舵机，channel 1 对应水平旋转舵机。PCA9685 的 I2C 地址为 0x40，All-Call 地址为 0x70。

关键接线如表~\ref{tab:wiring} 所示。实际接线表已整理为 Excel 文件，随项目源代码一并保存。

\begin{table}[htbp]
  \centering
  \caption{关键硬件接线}
  \label{tab:wiring}
  \begin{tabular}{lll}
    \toprule
    PCA9685/模块端口 & 连接对象 & 说明 \\
    \midrule
    VCC & OrangePi 40Pin 第 1 脚 3.3V & I2C 逻辑供电 \\
    GND & OrangePi 40Pin 第 6 脚 GND & 共地 \\
    SDA & OrangePi 40Pin 第 3 脚 SDA & I2C 数据线 \\
    SCL & OrangePi 40Pin 第 5 脚 SCL & I2C 时钟线 \\
    V+ / 绿色端子 & 独立 5V 舵机供电 & 舵机电源 \\
    channel 0 & 俯仰舵机 & 控制摄像头上下方向 \\
    channel 1 & 水平舵机 & 控制云台左右方向 \\
    USB 摄像头 & OrangePi USB 口 & 图像采集 \\
    \bottomrule
  \end{tabular}
\end{table}

\section{软件模块}

软件工程采用 Python 主程序与 C/C++ 扩展相结合的方式组织。主要模块包括：

\begin{itemize}
  \item \textbf{配置模块}：读取摄像头、视觉、控制、状态机、日志和硬件参数；
  \item \textbf{摄像头模块}：基于 OpenCV VideoCapture 获取图像帧；
  \item \textbf{视觉模块}：完成 HSV 阈值分割、形态学处理和轮廓筛选；
  \item \textbf{控制模块}：根据目标偏差计算 pan/tilt 角度；
  \item \textbf{硬件模块}：优先使用 smbus2 直接驱动 PCA9685，并保留 Adafruit 驱动与模拟硬件回退；
  \item \textbf{状态机模块}：管理 IDLE、TRACKING、LOST\_SHORT 和 SEARCH 状态；
  \item \textbf{日志模块}：记录帧率、检测耗时、控制耗时、偏差与舵机角度；
  \item \textbf{网页监控模块}：基于 HTTP/MJPEG 提供远程实时画面。
\end{itemize}

\section{控制策略}

控制模块以图像中心为期望位置，计算目标中心误差：

\begin{equation}
  e_x = c_x - \frac{W}{2}, \qquad e_y = c_y - \frac{H}{2}
\end{equation}

其中 $c_x,c_y$ 为目标中心坐标，$W,H$ 为图像宽高。控制器在误差较小时引入死区，避免舵机在目标接近中心时频繁抖动；在误差较大时输出角度增量，并通过最大步进、角度限位和平滑系数限制舵机动作幅度。当前配置中水平角范围为 20 到 160 度，俯仰角范围为 35 到 145 度，中心角均为 90 度。

